\section{Recreating the attack}
\label{sec:recreation}

\subsection{Design}

The lab is one Mininet network inside a Docker container. A supplier pipeline signs releases, an updater serves them, a DNS coordinator answers check-ins, and an HTTP receiver acts as the control server. The victim's only link to the switch crosses a bridging sensor, so every packet is captured rather than read from an application log. Five small web apps expose the supplier, the client, the operator, the presenter, and the defender.

The client verifies the catalog digest, an RSA/SHA-256 signature against a pinned public key, and the component digest before extraction. The supplier can publish V1 and V2 as data-only packages and one V3 package that declares a component. The manifest names \texttt{orion-improvement-program}, its file, and its SHA-256; the pipeline rejects files the manifest does not declare~\cite{task05lab}.

\subsection{Executed chain}

Table~\ref{tab:killchain} maps the incident to the lab. The status column distinguishes steps that ran from steps that exist only as an artifact.

\begin{table}[H]
\centering
\footnotesize
\caption{Kill-chain mapping. ``Represented'' means the lab contains the artifact but not the step that produced it.}
\label{tab:killchain}
\begin{tabular}{@{}>{\RaggedRight}p{2.3cm}>{\RaggedRight}p{4.7cm}>{\RaggedRight}p{4.7cm}>{\RaggedRight}p{2.5cm}@{}}
\toprule
\textbf{Phase} & \textbf{In the incident} & \textbf{In the lab} & \textbf{Status} \\
\midrule
Reconnaissance & Build chain understood before SUNSPOT was written~\cite{crowdstrike-sunspot} & Nothing & Not modelled \\
Weaponization & Source file replaced during the build, original restored~\cite{crowdstrike-sunspot} & Prebuilt V3 package declares a component and pins its digest & Represented \\
Delivery & Signed update through the normal channel~\cite{mandiant-campaign} & Signed package published to a stable channel, downloaded over HTTP & Executed \\
Exploitation & Customer installs the update & Client verifies and installs V3 & Executed \\
Installation & Backdoored DLL loads inside Orion & Installer registers the component and starts one process & Executed, reduced \\
Command and control & DGA check-in, mode by A record, CNAME handoff, OIP-shaped HTTPS & Encoded check-in, mode by coordinator answer, CNAME handoff, HTTP POST & Executed, reduced \\
Actions on objectives & TEARDROP to Cobalt Strike, AdFind, credential access, SAML forgery & One shell command and its output & Executed, reduced \\
\bottomrule
\end{tabular}
\end{table}

The client accepts the tampered package because every check passes. The catalog digest matches, the signature verifies against the pinned key, and the component digest matches the manifest. That is the property the lab was built to show: verification answers whether the supplier signed the artifact, not whether the build was clean.

After installation, the component waits three seconds, builds a base32 label from an XORed record holding a version byte, a mode byte, a user identifier, host and user names, and a CRC-16, and queries the coordinator. A bad checksum returns NXDOMAIN. The coordinator answers \texttt{192.0.2.10} for passive, \texttt{192.0.2.20} for active, or NXDOMAIN when disabled. The component starts at mode 4, moves to 5 on the activation address, and exits after a final mode-3 check-in when the coordinator is disabled.

Active traffic goes to \texttt{/swip/upd/orion.improvementprogram.asmx} with an \texttt{Orion/2020.2.4} agent string and a body of fields modelled on the Improvement Program: \texttt{customerId}, \texttt{machineName}, \texttt{userName}, \texttt{osVersion}, \texttt{ipAddress}, and \texttt{processId}. The receiver answers with \texttt{swip.tasks}, and the component returns output in a \texttt{telemetry} block. One shell command per task runs with \texttt{/bin/sh -lc}, with a guardrail against destructive patterns.

Evidence from the demo run is the client history for found, downloaded, verified, installed, and registered component; the capture at \texttt{/tmp/task05-poc/capture/c2-traffic.pcap}; and the sensor state at \texttt{/tmp/task05-poc/sensor/state.json}. The sample package and its digest are pinned in \texttt{pipeline/sample-v3-manifest.json}~\cite{task05lab}.

\subsection{What the recreation does not do}

The largest gap is weaponization. There is no build system, so SUNSPOT's replacement and restore is represented by a prebuilt package rather than performed. The lab has no Windows host, no signed PE, and no DLL loaded inside a product. Dormancy is compressed from 12 to 14 days to three seconds, and the environment checks are omitted, which makes the lab implant more portable than the original. Check-ins repeat every four seconds against a real minimum of one minute for HTTP and sleeps measured in hours for DNS. The control channel is plain HTTP so the capture can be read, and the implant queries the coordinator directly instead of using the system resolver. There is one victim, no persistence, no second stage, and no lateral movement.

These are stated here because Table~\ref{tab:killchain} makes the chain look more complete than the lab is. What the lab reproduces is the trust path and the shape of the check-in, not the actor's code.
